%% =====================================================================
%%  T-07-04  The Mirror
%%  -------------------------------------------------------------------
%%  One idea per file. Core is mandatory. Slots in canonical order.
%%  Max 700 words. No layout commands. No filler. New source material
%%  for this entry goes in sources/<BOOK>/T-07-04.tex -- never by
%%  rewriting this file (docs/RULES.md R0.1).
%% =====================================================================
%% @kind: topic
%% @id: T-07-04
%% @title: The Mirror
%% @subtitle: Repeat the last few words and silence does the rest
%% @chapter: c07
%% @order: 40
%% @status: stable
%% @sources: B5
%% @updated: 2026-09-19

\topic{T-07-04}{The Mirror}{Repeat the last few words and silence does the rest}

\begin{core}
Repeat the last one to three critical words the counterpart said, with upward tone, then stop. The similarity triggers an unconscious bonding response --- imitation is how humans comfort each other --- and the void you leave is filled by the counterpart elaborating, refining, and often correcting themselves toward the truth. It is collection with no questions asked.
\end{core}
\begin{mechanism}
Isopraxism: humans copy each other to comfort each other, in speech patterns, vocabulary, tempo and tone, mostly without noticing, as a sign of bonding and being in sync. A spoken mirror re-creates that signal on demand: the counterpart hears their own framing come back and reads it as deep understanding, so they keep talking to justify it. The silence after the mirror is half the tool --- people cannot tolerate a conversational void and plug it, and what they plug it with is unscripted. Because no question was asked, nothing was conceded and no direction was disclosed; the counterpart built the picture themselves and trusts it more for that.
\end{mechanism}
\begin{conditions}
Strongest when the counterpart is committed to their story and you need its structure, not its summary; when asking direct questions would signal your interest or your ignorance; and early in contact, before rapport exists to spend. Weakens badly when overused --- three mirrors in a row reads as mockery or obtuseness --- and when the counterpart is themselves reading you carefully, in which case the technique is visible and can be named.
\end{conditions}
\begin{application}
  \item Let them finish a substantive statement. Do not nod through it; actually hold the last clause in memory.
  \item Repeat the critical one to three words back with a slight upward tone --- the tone of a question you are not quite asking.
  \item Say nothing for two to four seconds. Let the silence do the pull.
  \item Collect what comes back. The elaboration is the product; the mirror was only the pump.
  \item When the flow stalls, mirror the new critical phrase and repeat the cycle. Two to four cycles is a working ceiling.
\end{application}
\begin{feedback}
  \item Landing: the counterpart expands, softens, or corrects their own statement without being challenged.
  \item Landing: tempo drops and personal details start appearing --- a sign the guard is down.
  \item Failing: the counterpart answers your mirror as a question, briefly, and goes quiet themselves.
  \item Failing: they ask what you mean, or mirror you back. The tool has been seen.
\end{feedback}
\begin{failure}
Overuse converts rapport-building into an obvious trick, and the counterpart recalibrates: if you were faking understanding, what else was performed? A mirror aimed at the wrong word --- an incident, not the stake --- signals that you were not listening at all, which costs more than not mirroring. Against someone trained or naturally guarded, mirrors read as interrogation rhythm and produce clipped, managed answers.
\end{failure}
\begin{countermeasures}
  \item Notice a repeated phrase coming back at you with question-tone; answer the implied question once, then ask a direct one of your own. Mirrors cannot survive direct questions.
  \item Treat your own words in someone else's mouth as a prompt to become briefer, not more elaborate.
  \item Use silence awareness deliberately: when you feel the pull to fill a pause, the pause is working on you --- make the other side wait too.
  \item Slow, quiet speech from a counterpart mid-confrontation is a de-escalation grip, not timidity; re-assess before you relax.
\end{countermeasures}

\sources{B5}
\seesources
\seealso{T-07-01, T-07-02, T-13-02}
